\SourceLeaf{192}eigentliche Lösung des Gordischen Knotens. Allerdings kann man zugeben, dass diese Auslassung
hier insofern wenig ins Gewicht fällt, als sich die Beweisführung des Verfassers zur Unmöglichkeit des
dritten Falls, wenn sie zulässig wäre, unmittelbar auch auf diesen Fall erstrecken ließe.
Doch gerade diesem angeblichen Beweis der Unmöglichkeit des dritten Falls können wir keine
Gültigkeit zugestehen. Der Verfasser stellt die Sache so dar. Wenn man zum Beispiel im dritten Fall
annimmt, dass die ersten beiden Stücke gleich sind und das dritte größer ist, wäre
$DC$ also größer als $CL$. Da aber $AML$ ebenfalls ein gleichschenkliges Dreieck ist,
auf das dieselbe Grundbedingung wie auf das ursprüngliche Dreieck $AON$ zutrifft, müsste, wenn
der dritte Fall mit seiner angenommenen Unterteilung zuträfe, $DC=CL$ gelten,
im Widerspruch zum zuvor Gefundenen.

Wir glauben, an dieser Stelle des Beweises das, worauf es nach Ansicht des Verfassers ankommt, noch etwas
klarer und bestimmter angedeutet zu haben als er selbst,
wodurch sich aber auch die Schwäche des Beweises, wie uns scheint, leichter erkennen lässt. Denn
offenbar wird hier ganz willkürlich angenommen, dass bei allen gleichschenkligen Dreiecken mit
dem Winkel $A$ an der Spitze und einem größeren Winkel an der Basis, wenn man die im Lehrsatz
angegebene Konstruktion an ihnen vornimmt, die Folge der abgeschnittenen Stücke hinsichtlich
ihres Gleichbleibens, Größer- oder Kleinerwerdens in jedem Fall, unabhängig von der Länge
der Seiten, notwendig dieselbe sein muss. Eine solche Annahme kann unmöglich als selbstverständlich
gelten. Da sich aber der versuchte Beweis der
Unmöglichkeit des dritten (wie auch des vierten und fünften) Falls allein darauf stützt und der ganze Artikel
auch keine anderen Mittel bietet, um die Unmöglichkeit des übergangenen Falls zu beweisen,
glauben wir, damit das oben ausgesprochene Urteil hinreichend gerechtfertigt zu haben.
Dabei lassen wir der ganzen übrigen einfallsreichen Durchführung in den folgenden
Artikeln gern volle Gerechtigkeit widerfahren.

\begin{center}\rule{3cm}{0.4pt}\end{center}

\Needspace{10\baselineskip}
\begin{center}\textsc{Bemerkung.}\end{center}

{\small
Diese bereits in Band IV, Seiten 368 bis 370 abgedruckte Anzeige wurde hier der Vollständigkeit
wegen wiedergegeben. Die Figur, die sich weder im Original noch im ersten
Abdruck findet, soll das Verständnis des Textes erleichtern.

\begin{flushright}\textsc{Stäckel.}\end{flushright}
}

\begin{center}\rule{3cm}{0.4pt}\end{center}
